\section{Limitations}
\label{sec:limitations}

\paragraph{Access to the commercial model.} Jev~1.13 is available only through
its HTTP API. Its architecture, parameter count, tokenizer, training data and
serving configuration are unpublished. The vendor describes its training as
RLCD, Reinforcement Learning for Calibrated Decisions, and states that it
optimises for ``epistemically honest probabilities'' \citep{typesafe2026jev},
but no paper, objective function or reward definition has been released, so we
attribute no measured behaviour to a training method. The returned model string
is the only version identifier, which means server-side changes between our run
and a replication cannot be ruled out, and whether any slice appears in its
training data cannot be determined. Our results for Jev therefore describe the
deployed service as a customer reaches it.

\paragraph{D1 access and incomplete coverage.} The D1 alias \texttt{d1:free}
does not identify an immutable checkpoint. We cannot separate model changes
from service changes or infer its hardware, training exposure or self-hosting
licence from the API response. Two rate-limited MMLU-Pro cases remain absent
from the results reported here. D1 comparisons use shared successful cases and
do not include paired confidence intervals in this revision. Its client
concurrency changed during the run, so its latency is not a controlled speed
comparison with the other systems.

\paragraph{Evaluation design.} We wrote the instructions and criteria for every
slice ourselves; they are listed in Appendix~\ref{app:prompts}, and different
wording could shift the results. Each request carries one question, so the
batched setting in which this model class is reported to be fastest is not
measured. The four systems also run on different hardware: Jev and D1 remotely, so their
latency includes a network round trip, djev on one A100, and Laya on a laptop.
Latency figures therefore compare systems, not models. Each system was run once.
Language-model inference is not deterministic by default, since GPU kernel
results depend on how many requests are batched together
\citep{he2025nondeterminism}; djev ran with a fixed seed but six requests at a
time, and the hosted API offers no seed control. Small differences between
systems should be read with this in mind.

\paragraph{Data.} Five slices are drawn from datasets that appear, or are likely
to appear, in the training data of these models
(Appendix~\ref{app:provenance}), and results should not be averaged across
exposed and unexposed slices. Most safety slices are sampled close to class balance, so their
accuracies do not estimate incident rates in deployment. Two sources do not
state their label polarity; the scorer reports both directions as a diagnostic.
A higher score under one direction does not verify its semantics. These source
mappings require independent verification before substantive safety conclusions.
The harness caps states
at 12{,}000 characters, which truncated 34 ATBench cases and one jailbreak case.

\paragraph{Laya's context length.} Laya's default context is 512 tokens, and it
truncated every ATBench trajectory internally; our truncation flag, which tracks
only the harness's own character cap, did not detect this. ATBench trajectories
average about 1{,}500 tokens, so Laya's 65.4\% on this slice
(Table~\ref{tab:main}) reflects a truncated input and is not directly comparable
with Jev and djev, which received substantially longer inputs
(Table~\ref{tab:context}).
